\CourseSection{Conclusion}\label{sec:conclusion}
\Guidance{Summarize the project and the evidence established in Sections 1--5. End with the expected contribution and the team's readiness to proceed.}

This project aims to examine size, churn, age, clone coverage, and readability within AI agent skills using the GitSkills dataset. The team plans to study how these skills are developed and maintained, and whether traditional software engineering metrics can be meaningfully applied to them. The final system will be a Python analysis pipeline that processes the dataset and calculates metrics based on the available fields.

Development will follow six milestones (M1–M6). M1 focuses on dataset preparation and verification of sample records. Next, M2 involves the initial metric calculations and checking them against reference examples. Once the methods have been established, the team will calculate the supported metrics across the dataset and generate preliminary tables and visualizations (M3). This is followed by M4, which involves reviewing results and documenting focused research questions. Based on these findings, the team will investigate follow-up questions and potential extensions during M5. Finally, the pipeline will undergo validation and reproducibility testing in M6. The pipeline will be designed to process the full dataset within 24 hours while using no more than 16 GB of memory. Testing will include unit tests, integration tests, and reproducibility checks to ensure the calculated metrics are accurate and consistent. Since the project relies on publicly available data and open-source tools, the estimated total cost is \$0 CAD.

The expected outcome of the project is to identify and document meaningful metrics that can help developers and researchers better understand AI agent skills and how they evolve. The full scope of the project will depend on the results of preliminary analysis, as focused research questions will be established during M4 and investigated further during M5. Although this makes part of the project's scope open-ended, the team has established a clear development plan, assigned milestones, and identified the resources needed for success. With initial pilot analyses already completed, the team is prepared to move forward with implementation and validation.
